% paper/sections/appendix_legacy_methods.tex
%
% 付録 F: 標準手法の詳細（CSF モデル・Rhie--Chow 補正）
% 本付録は本稿が GFM + DCCD 定式化に移行する以前の
% CSF および Rhie--Chow 手法の完全な定式化を参考のため収録する．
%

\section{標準手法の詳細：CSF モデルと Rhie--Chow 補正}
\label{app:legacy_methods}

% ═══════════════════════════════════════════════════════════════════════════════
\subsection{CSF（連続表面力）モデルの完全な定式化}
\label{app:csf_full}
% ═══════════════════════════════════════════════════════════════════════════════

Navier-Stokes 方程式に表面張力項 $\bm{f}_\sigma^*$ を組み込む前に，
その導出を示す必要がある．表面張力は本来「界面上の線力（3次元では面力）」だが，
これをそのまま体積方程式に入れることはできない．
\textbf{連続表面力（CSF: Continuum Surface Force）モデル} \cite{Brackbill1992} は
この問題を解決する標準的手法である．

% ─────────────────────────────────────────────────────────────────────────────
\subsubsection{なぜデルタ関数が必要か：面積分から体積積分への変換}
% ─────────────────────────────────────────────────────────────────────────────

\paragraph{直感的理解：「フィルター」としてのデルタ関数}
表面張力は本来，界面 $\Gamma$ という\textbf{面（2次元曲面）}の上にのみ作用する力である．
一方，Navier-Stokes 方程式は\textbf{体積（3次元領域）}全体にわたる方程式である．
「面の上の力」を「体積全体の力」として書き直すには，
\textbf{「界面でのみ1，それ以外でゼロ」という数学的なフィルター}が必要になる．
これがディラックデルタ関数 $\delta_s$ の役割である．

具体的には，表面張力による力 $F_\sigma$ を体積積分と面積分の両方で表すと：
\[
  \underbrace{\int_\Omega \bm{f}_\sigma(\bm{x})\, dV}_{\text{体積積分（NS 方程式側）}}
  = \underbrace{\int_\Gamma \sigma\kappa\hat{\bm{n}}\, dS}_{\text{面積分（物理的表現）}}
\]
この等式が成り立つためには，$\bm{f}_\sigma = \sigma\kappa\hat{\bm{n}}\,\delta_s(\bm{x})$
（界面上では $+\infty$，それ以外では $0$，積分すると界面の面積を与える関数）
とおく必要がある．

\paragraph{1次元での直感：ヘヴィサイド関数の微分}
\begin{figure}[htbp]
\centering
\begin{tikzpicture}[scale=1.0]
  %% 左：ヘヴィサイド関数
  \begin{scope}[xshift=-45mm]
    \node[font=\small\bfseries, text=navy] at (0, 2.8) {（a）ヘヴィサイド関数 $H_\varepsilon(\phi)$};
    \draw[->] (-2.2,0) -- (2.2,0) node[right, font=\small]{$\phi$};
    \draw[->] (0,-0.3) -- (0, 2.5) node[above, font=\small]{$H_\varepsilon$};
    \draw[very thick, blue2, domain=-2:2, samples=80]
      plot (\x, {1.8/(1+exp(-3*\x))});
    \node[font=\scriptsize, blue2] at (1.4, 1.65) {$H_\varepsilon(\phi)$};
    \draw[dashed, gray] (-2.2, 1.8) -- (2.2, 1.8) node[right, font=\scriptsize, gray]{$1$};
    \draw[dashed, gray] (-2.2, 0)   -- (2.2, 0)   node[right, font=\scriptsize, gray]{$0$};
    \node[font=\scriptsize, red2] at (0.15, -0.25) {$\phi{=}0$};
    \node[font=\scriptsize] at (-1.3, 0.9) {液相};
    \node[font=\scriptsize] at (1.3,  0.9) {気相};
  \end{scope}
  %% 右：デルタ関数（微分）
  \begin{scope}[xshift=45mm]
    \node[font=\small\bfseries, text=navy] at (0, 2.8) {（b）微分 $\delta_\varepsilon(\phi) = H_\varepsilon'(\phi)$};
    \draw[->] (-2.2,0) -- (2.2,0) node[right, font=\small]{$\phi$};
    \draw[->] (0,-0.3) -- (0, 2.5) node[above, font=\small]{$\delta_\varepsilon$};
    \draw[very thick, red2, domain=-2:2, samples=80]
      plot (\x, {2.0*exp(-3*\x*\x)});
    \fill[red2!15, domain=-2:2, samples=80]
      plot (\x, {2.0*exp(-3*\x*\x)}) -- (2, 0) -- (-2, 0) -- cycle;
    \node[font=\scriptsize, red2] at (0.6, 1.8) {$\delta_\varepsilon(\phi)$};
    \node[font=\scriptsize, red2] at (0.15, -0.25) {$\phi{=}0$};
    \draw[<->, thick, darkgray]
      (-0.6, -0.15) -- (0.6, -0.15) node[midway, below, font=\scriptsize]{幅 $\approx 2\varepsilon$};
    \node[font=\scriptsize, darkgray] at (1.6, 0.3) {界面外 $\approx 0$};
    \node[font=\scriptsize] at (0, 1.0) {$\int\delta_\varepsilon\,d\phi = 1$};
  \end{scope}
  %% 微分矢印
  \draw[-Stealth, very thick, navy] (-0.8, 1.4) -- (0.8, 1.4)
        node[midway, above, font=\small, navy]{$\dfrac{d}{d\phi}$};
\end{tikzpicture}
\caption{（a）CLS 変数 $\psi = H_\varepsilon(\phi)$（ヘヴィサイド関数の滑らか版）．
界面 $\phi=0$ を中心に 0 から 1 へ連続的に変化する．
（b）その微分 $\delta_\varepsilon(\phi) = H_\varepsilon'(\phi)$ は界面近傍にのみ集中した
ピーク（デルタ関数の滑らか版）を形成する．
$|\bnabla\psi| = \delta_\varepsilon(\phi)|\bnabla\phi| \approx \delta_\varepsilon(\phi)$
が界面デルタ関数 $\delta_s$ の数値近似となる．}
\label{fig:heaviside_delta_app}
\end{figure}

1次元の場合，ヘヴィサイド関数 $H(x)$（$x>0$ で 1，$x<0$ で 0）の微分は
ディラックデルタ $\delta(x)$ である：
\[
  \frac{d}{dx}H(x) = \delta(x)
\]
数値計算では厳密なステップ関数の代わりに，
滑らかな近似 $H_\varepsilon(\phi)$（図\ref{fig:heaviside_delta_app}(a)）を用いる．
その微分 $H_\varepsilon'(\phi) = \delta_\varepsilon(\phi)$ は
「幅 $\approx 2\varepsilon$ を持つ滑らかなピーク関数」（図\ref{fig:heaviside_delta_app}(b)）となり，
これが界面ディラックデルタ $\delta_s$ の数値的な近似として機能する．

% ─────────────────────────────────────────────────────────────────────────────
\subsubsection{物理的表面張力と界面のディラックデルタ}
% ─────────────────────────────────────────────────────────────────────────────

界面 $\Gamma$ 上に作用する表面張力は，単位面積あたり $\sigma\kappa\hat{\bm{n}}$ の法線力である．
これを全空間で定義された体積力 $\bm{f}_\sigma$ として書くには，
界面上だけ値を持つ\textbf{界面ディラックデルタ} $\delta_s$ を用いる：
\begin{equation}
  \bm{f}_\sigma = \sigma\kappa\hat{\bm{n}}\,\delta_s(\bm{x})
  \label{eq:csf_dirac_app}
\end{equation}
$\delta_s$ は界面 $\Gamma$ 上では $+\infty$，それ以外では $0$ の一般化関数であり，
$\int_\Omega \delta_s(\bm{x})\,dV = \int_\Gamma dS$（面積分）という性質を持つ．

% ─────────────────────────────────────────────────────────────────────────────
\subsubsection{CLS を用いた \texorpdfstring{$\delta_s$ と $\hat{\bm{n}}$}{δ\_s と n-hat} の近似}
% ─────────────────────────────────────────────────────────────────────────────

解析的には $\delta_s$ は扱いにくいが，CLS 変数 $\psi = H_\varepsilon(\phi)$ を用いると
\textbf{滑らかな近似}が得られる：
\begin{align}
  |\bnabla\psi| &= |H_\varepsilon'(\phi)|\,|\bnabla\phi|
                 = \delta_\varepsilon(\phi)\,|\bnabla\phi|
                 \approx \delta_\varepsilon(\phi)
  \label{eq:csf_delta_psi_app}
\end{align}
最後の近似は $|\bnabla\phi|\approx 1$（符号付き距離関数の Eikonal 性質）による．
すなわち，\textbf{$|\bnabla\psi|$ が界面ディラックデルタ $\delta_s$ の滑らかな近似}となる．

以下の精度証明は $H_\varepsilon$ の具体的な形式に依拠する（正式導出は第\ref{sec:levelset}章）．
シグモイド型滑らかヘヴィサイド関数を
\begin{equation}
  H_\varepsilon(\phi) \equiv \frac{1}{1+e^{-\phi/\varepsilon}}
  \label{eq:Heps_def_app}
\end{equation}
と定義する（$\phi=0$ で $H_\varepsilon = 1/2$，$\phi\to-\infty$ で $0$，$\phi\to+\infty$ で $1$）．
その微分は閉形式で
\begin{equation}
  \delta_\varepsilon(\phi) \equiv H_\varepsilon'(\phi)
  = \frac{1}{\varepsilon}H_\varepsilon(\phi)\bigl(1-H_\varepsilon(\phi)\bigr)
  \label{eq:delta_eps_def_app}
\end{equation}
と書け，かつ $\delta_\varepsilon(\phi) = \delta_\varepsilon(-\phi)$（偶関数性）が成り立つ．

\textbf{近似精度（付録~\ref{app:csf_delta_precision} に証明）：}
Eikonal 条件 $|\bnabla\phi|=1$ のもとで，法線座標 $s$ のテイラー展開と
$\delta_\varepsilon$ の偶関数性から：
\[
  \int_\Omega f\,|\bnabla\psi|\,dV - \int_\Gamma f\,dA_\Gamma
  = \Ord{\varepsilon^2\,\kappa^2\,\|f\|_{L^\infty}}
\]
$\varepsilon \sim \Delta x$ の設定では誤差は $\Ord{\Delta x^2}$ である．

一方，界面の外向き単位法線は：
\begin{equation}
  \hat{\bm{n}} = \frac{\bnabla\psi}{|\bnabla\psi|}
  \label{eq:csf_normal_app}
\end{equation}

% ─────────────────────────────────────────────────────────────────────────────
\subsubsection{CSF 体積力の最終形：\texorpdfstring{$\bm{f}_\sigma = \sigma\kappa\bnabla\psi$}{f\_sigma = sigma*kappa*grad(psi)}}
% ─────────────────────────────────────────────────────────────────────────────

式 \eqref{eq:csf_delta_psi_app} と \eqref{eq:csf_normal_app} を式 \eqref{eq:csf_dirac_app} に代入する：
\begin{equation}
  \bm{f}_\sigma = \sigma\kappa\,\hat{\bm{n}}\,\delta_s
  \approx \sigma\kappa\,\frac{\bnabla\psi}{|\bnabla\psi|}\,|\bnabla\psi|
  \approx \sigma\kappa\,\bnabla\psi
  \label{eq:csf_final}
\end{equation}

ただし，この近似は $|\bnabla\phi| \approx 1$（Eikonal 条件）が成立する場合に限り有効である．
CLS 法では再初期化ステップによって $|\bnabla\phi|\approx 1$ を維持し，この近似の精度を保証する．

$\bnabla\psi$ は界面の法線方向を向き，界面で最大値をとる．
したがって $\sigma\kappa\bnabla\psi$ は界面近傍でのみ作用し，
全空間積分すると Young-Laplace のジャンプ条件 $[p]_\Gamma = \sigma\kappa$ を再現する
（以下の検証を参照）．

\medskip
\noindent\textbf{検証：CSF モデルが Young-Laplace ジャンプ条件を再現すること}
\label{sec:csf_young_laplace_app}

CSF 体積力 $\bm{f}_\sigma = \sigma\kappa\bnabla\psi$ が界面ジャンプ条件
$[p]_\Gamma \equiv p_{\text{gas}} - p_{\text{liq}} = \sigma\kappa$
を再現することを，界面法線方向の積分によって示す．

平衡状態（速度ゼロ）では運動方程式より
$\bnabla p = \sigma\kappa\bnabla\psi$ が成立する．
界面を法線方向 $s$ に沿って積分すると：
\[
  p_{\text{gas}} - p_{\text{liq}}
  = \sigma\kappa \int_{-\infty}^{+\infty} \frac{\partial\psi}{\partial s}\,ds
  = \sigma\kappa\,(1 - 0) = \sigma\kappa
\]
\begin{equation}
  p_{\text{gas}} - p_{\text{liq}} = \sigma\kappa
  \label{eq:young-laplace-app}
\end{equation}

% ═══════════════════════════════════════════════════════════════════════════════
\subsection{Rhie--Chow 補正の完全な定式化}
\label{app:rhie_chow_full}
% ═══════════════════════════════════════════════════════════════════════════════

本付録では，第\ref{sec:collocate}章§\ref{sec:checkerboard_solution} で参照した
Rhie--Chow 補正の完全な定式化を収録する．

\subsubsection{補正された界面速度の定義}
本稿では，Chung~\cite{Chung2002} に基づく変密度場に対応した形式を用いる．
東側フェイス $e$ における $x$ 方向速度 $u_e$ は，隣接するセル中心 $P$ と $E$ の値を用いて以下のように定義される．
\begin{equation}
 u_e = \overline{u_e} - d_e \left[ \left(\pd{p}{x}\right)_e - \overline{\left(\pd{p}{x}\right)_e} \right]
 \label{eq:rc-concept-app}
\end{equation}
ここで，
\begin{itemize}
    \item $\overline{u_e}$: セル中心 $P, E$ の速度を線形補間した速度．$\overline{u_e} = (u_P+u_E)/2$．
    \item $(\partial p/\partial x)_e$: フェイス $e$ での圧力勾配．単純な中心差分 $\frac{p_E-p_P}{\Delta x}$ で計算．
    \item $\overline{(\partial p/\partial x)_e}$: セル中心 $P, E$ で評価された圧力勾配を線形補間したもの．
    $\overline{(\partial p/\partial x)_e} = \frac{1}{2}\left[ \left(\frac{\partial p}{\partial x}\right)_P + \left(\frac{\partial p}{\partial x}\right)_E \right]$
    \item $d_e$: 補間係数．運動量方程式から導かれ，時間刻み幅 $\Delta t$ と密度 $\rho$ に依存する．
    $d_e = (\Delta t / \rho)_e$（式~\eqref{eq:rc-face} 参照）．
\end{itemize}

Rhie--Chow 式 \eqref{eq:rc-concept-app} に2種類の補間が現れる：
\begin{itemize}
  \item \textbf{$\bar{u}_f$（速度）→ 算術平均}：
    速度は気液界面をまたいでも物理的に連続であるため，単純な線形補間が適切である．
  \item \textbf{$d_f = \Delta t (1/\rho)_f^{\mathrm{harm}}$（密度係数）→ $1/\rho$ の調和平均}：
    密度 $\rho$ は界面で不連続であり，フラックスの連続性（直列抵抗則）から
    $a \equiv 1/\rho$ の調和平均
    $(1/\rho)_f^{\mathrm{harm}} = 2a_La_R/(a_L+a_R) = 2/(\rho_L+\rho_R)$
    が自然に導かれる（第\ref{sec:pressure}章；
    付録~\ref{app:fvm_face_coeff} の FVM 積分論理を参照）．
    これは $\rho$ の算術平均 $(\rho_L+\rho_R)/2$ の逆数に等しい．
    注：$\rho$ の調和平均の逆数 $(\rho_L+\rho_R)/(2\rho_L\rho_R)$ とは異なる
    （付録~\ref{app:fvm_face_coeff} の「よくある実装ミス」を参照）．
\end{itemize}

\paragraph{補正項の働き}
式 \eqref{eq:rc-concept-app} の第2項が Rhie--Chow 補正の核心である．
チェッカーボード圧力場では，フェイスでの勾配 $(p_E-p_P)/\Delta x$ は大きな値を持つが，
セル中心勾配の平均はゼロに近くなるため，この差が非物理的な振動を効果的に減衰させる．

\subsubsection{本研究における離散化形式}
\label{sec:rc_implementation}
最終的に，圧力 Poisson 方程式との整合性を厳密に保つため，本研究では以下の離散化形式を用いる
（表面張力補正項を加えた完全形式 \eqref{eq:rc-face-balanced} については §\ref{sec:rc_balanced_force} を参照）．
\begin{equation}
 u_f = \bar{u}_f - \Delta t\left( \frac{1}{\rho^{n+1}} \right)^{\mathrm{harm}}_f \left[ (\bnabla p^n)_f - \overline{(\bnabla p^n)}_f \right]
 \label{eq:rc-face}
\end{equation}
ここで，$(\nabla p^n)_f$ はフェイスにおける前時刻圧力 $p^n$ の勾配（陽的評価），
$\overline{(\nabla p^n)}_f$ はセル中心における前時刻圧力勾配の補間値である．
係数 $(1/\rho^{n+1})_f^\mathrm{harm} = 2/(\rho^{n+1}_P + \rho^{n+1}_E)$ には現時刻の $1/\rho$ の調和平均（$= \rho$ の算術平均の逆数）を用いる．

\paragraph{補間方式の根拠}
$1/\rho$ の算術平均と調和平均の差は密度比が大きいほど顕著になる：
\begin{center}
\begin{tabular}{lccc}
\hline
流体系 & $\rho_l/\rho_g$ & $(1/\rho)^{\mathrm{arith}}_f \cdot \rho_g$ & $(1/\rho)^{\mathrm{harm}}_f \cdot \rho_g$ \\
\hline
水・水蒸気 & $10$ & $0.55$ & $0.18$ \\
水・空気（常温） & $1000$ & $\approx 0.500$ & $\approx 0.002$ \\
\hline
\end{tabular}
\end{center}

\subsubsection{Rhie--Chow 補正済み発散の定義}
\label{sec:rc_divergence}

Rhie--Chow 補正済みフェイス速度を用いた\textbf{フェイス速度の発散}：
\begin{equation}
  \bigl(\bnabla_h^{\mathrm{RC}}\!\cdot\bu^*\bigr)_{i,j}
  \equiv
  \frac{(u_f)_{i+1/2,j} - (u_f)_{i-1/2,j}}{\Delta x}
  + \frac{(v_f)_{i,j+1/2} - (v_f)_{i,j-1/2}}{\Delta y}
  \label{eq:rc_divergence}
\end{equation}
ここで $(u_f)_{i+1/2,j}$ は式~\eqref{eq:rc-concept-app} の Rhie--Chow 補正済みフェイス速度．
$p^n$ は前時刻圧力（陽的補正），$(1/\rho)^{\mathrm{harm}}_{i+1/2,j} = 2/(\rho^{n+1}_{i,j}+\rho^{n+1}_{i+1,j})$
は時刻 $n+1$ の密度値から評価する．

\noindent\textbf{【重要】PPE 右辺での Rhie--Chow 発散の必須使用：}
PPE 右辺の発散計算には，必ずこのフェイス速度の発散 $\nabla_h^{\mathrm{RC}}\cdot\bu^*$
を使用しなければならない
（セル中心発散の代用はチェッカーボードの逆流を招く）．

\noindent\textbf{注：本稿（GFM+DCCD 定式化）では}
DCCD フィルタ（§\ref{sec:dccd_decoupling}）+ GFM（§\ref{sec:gfm}）が
Rhie--Chow 補正を代替し，式~\eqref{eq:gfm_ccd_div} が PPE 右辺を構成する．

\subsubsection{Rhie--Chow 補正への表面張力項の拡張：フェイス速度への Balanced-Force 反映}
\label{app:rc_balanced_force}
\label{app:csf_balanced_force}

Balanced-Force の論理をフェイス速度の計算にも完全に反映させるためには，
表面張力体積力の差分項を加えた拡張形式が必要である：
\begin{equation}
  u_f = \bar{u}_f
    - \Delta t\left(\frac{1}{\rho^{n+1}}\right)_f^{\mathrm{harm}}
    \Biggl[
      \Bigl((\bnabla p^n)_f - \overline{(\bnabla p^n)}_f\Bigr)
      - \Bigl((\bm{f}_\sigma^{n+1})_f - \overline{(\bm{f}_\sigma^{n+1})}_f\Bigr)
    \Biggr]
  \label{eq:rc-face-balanced-app}
\end{equation}

ここで各量の定義は以下の通りである：
\begin{itemize}
  \item $(\bm{f}_\sigma^{n+1})_f = (\kappa_f^{n+1}/We)\cdot(\psi_E^{n+1}-\psi_P^{n+1})/h$：
        フェイス $f$ での表面張力体積力．
  \item $\overline{(\bm{f}_\sigma^{n+1})}_f = \tfrac{1}{2}[(\bm{f}_\sigma^{n+1})_P + (\bm{f}_\sigma^{n+1})_E]$：
        セル中心値の算術平均補間値．\textbf{必須：} セル中心値には CCD 演算子を使用する
        （Balanced-Force 整合性条件，§\ref{sec:balanced_force}）．
\end{itemize}

平衡状態では，フェイス直接差分レベルで $(\bnabla p^n)_f \approx (\bm{f}_\sigma^{n+1})_f$
がキャンセルし，セル中心平均で同一の CCD 演算子を使うことで
$\overline{(\bnabla p^n)}_f - \overline{(\bm{f}_\sigma^{n+1})}_f \to 0$ となる
（詳細は付録~\ref{app:balanced_force_taylor}）．
